\section{Términos clave: índice de palabras clave de este episodio}

Este episodio tiene una densidad de términos muy alta, y muchos de ellos cambian de significado según el contexto. La siguiente tabla reúne los términos centrales para ayudar al lector a repasar.

\begin{longtable}{p{0.23\textwidth}p{0.34\textwidth}p{0.34\textwidth}}
\toprule
Término & Explicación en una frase & Papel en este episodio \\
\midrule
AI War & La competencia en IA como una carrera armamentista que nadie puede permitirse perder & Sustituye el relato de una simple burbuja y explica por qué los gigantes tecnológicos y los Estados siguen invirtiendo. \\
AI Bubble & Debate sobre si las valuaciones y el gasto de capital están sobrecalentados & Pregunta de apertura del programa, pero no el marco de análisis definitivo. \\
SOTA & State of the art, el nivel más avanzado actual & Se advierte que no basta con mirar los benchmarks: también hay que considerar usuarios e ingresos. \\
Benchmark & Tabla de clasificación o métrica que mide la capacidad de un modelo con un conjunto de prueba & Tiene valor de referencia, pero sus puntuaciones pueden inflarse artificialmente o alejarse de las tareas reales. \\
ARR & Annual Recurring Revenue, ingresos recurrentes anuales & Mide la velocidad de comercialización de las empresas AI Native. \\
DAU/MAU & Usuarios activos diarios/mensuales & Compara la frecuencia de uso de ChatGPT y Gemini y su lugar en la mente de los usuarios como punto de entrada. \\
Preentrenamiento & Entrenar un modelo base con cantidades masivas de datos & Se compara con el petróleo: abundante, pero cada vez más cerca de su límite. \\
RL/Post-training & Seguir moldeando el modelo con retroalimentación, preferencias y trayectorias de expertos & Se compara con las energías nuevas: útil, pero limitado en volumen total y en costo. \\
Online Learning & Aprender de la interacción continua y mejorar el comportamiento futuro & Tercer paradigma de este episodio y variable decisiva para cambiar el panorama. \\
Continuous Learning & Aprendizaje continuo; suele usarse como casi sinónimo de Online Learning & Pone énfasis en que el modelo absorba nuevas experiencias a largo plazo. \\
Memory & Capacidad de conservar el contexto del usuario, de las tareas y del historial & Es una de las infraestructuras base de los Agents proactivos y del Online Learning. \\
LoRA & Low-Rank Adaptation, adaptación eficiente del modelo mediante parámetros de rango bajo & El programa la menciona como una de las infraestructuras del nuevo paradigma. \\
MoE & Mixture of Experts, activa solo una parte de las redes expertas según la entrada & Explica que los parámetros activados del modelo no necesariamente crecen sin límite, mientras aumenta la dispersión. \\
Agentic Web & Al convertirse el Agent en punto de entrada, se reorganizan las páginas web, las herramientas, los pagos y la distribución de las plataformas & Afecta la publicidad en buscadores, las Super Apps y las comisiones por transacción. \\
Modelo como producto & La experiencia de producto que percibe el usuario depende cada vez más de la capacidad del modelo & Muestra que una mejora del modelo cambia directamente el valor de las aplicaciones. \\
Datos como modelo & Las diferencias entre modelos provienen, en última instancia, de la distribución de los datos y del ciclo cerrado de retroalimentación & Muestra que las empresas emergentes deben contar con datos diferenciados. \\
Destilación de expertos & Convertir el proceso de trabajo de expertos de alto nivel en datos de entrenamiento & Elemento decisivo para extender horizontalmente las capacidades a distintas industrias. \\
L3/L4 local & Analogía con los niveles de la conducción autónoma para describir una automatización alta en escenarios específicos & Un objetivo comercial más realista que una AGI integral. \\
Neo Labs & Grupo de nuevos laboratorios surgidos de grandes empresas y de OpenAI & Exploran el siguiente paradigma, los sistemas multiagente, los materiales y los modelos de código abierto. \\
Robotics & Inteligencia robótica y sistemas de ejecución física & Dirección del mundo real que pone a prueba la generalización y la coordinación entre datos y hardware. \\
\bottomrule
\end{longtable}

\subsection{Resumen de la sección}

Todos estos términos apuntan a lo mismo: la competencia en IA ya no se limita a la capacidad de un modelo aislado, sino que se ha ampliado a un sistema compuesto de datos, producto, puntos de entrada, capacidad de cómputo, organización, capital y estrategia nacional. Entender cómo se relacionan los términos importa más que memorizar la clasificación de un trimestre.
